\begin{acknowledgement}
行文至此，本文的撰写工作已经接近尾声，我的本科学习生涯也即将画上一个句号。在此，我想借此机会向所有在这段旅程中给予我支持和帮助的人表达我最诚挚的感谢。

首先，我要感谢华东师范大学。这所学校不仅提供了优质的教育资源和学术环境，更是培养了我独立思考和创新能力的摇篮。在这里，我遇到了许多优秀的老师和同学，他们的智慧和热情激励着我不断前行。在这四年的学习期间，华师大的人文素养也深深地感染了我，一些极为优秀的通识课程让我收获良多，以至于时常回忆起这些课程时，不免感叹当时没有好好珍惜。

我要特别感谢我的导师，苏亭教授。
这篇论文的完成离不开苏教授的悉心指导。
苏教授不仅在学术上给予了我无私的指导和支持，还时常关心我的成长和发展。他严谨治学态度和对学生的关怀让我深受感动，也激励着我不断追求卓越，非常期待在未来的学习和研究中继续得到他的指导和帮助。

我要感谢王可以学姐。她在这篇论文的撰写过程中给予了我极大的帮助和支持。在研究方法与实验评估的设计以及最后的论文撰写过程中，学姐都提供了宝贵的建议和指导，并提出了很多修改意见。她的专业知识和耐心帮助我克服了许多困难，使得这篇论文能够顺利完成。

感谢符俊哲学长和蔡耀德学长，他们在方法设计和实验上提供了很多帮助，与他们的讨论和交流让我受益很多；感谢周炽金老师、孙海英老师和万成城老师。周老师、孙老师和万老师每次都来参加我们的讨论，为我提供了很多珍贵的建议；感谢上海工业控制安全创新科技有限公司的张绍景和柳泽，与他们的交流为我提供了一些新的思路。

在毕业论文之外，我要感谢的人还有很多——

感谢可爱的同学们，与你们的相处让我在大学生活中收获了许多美好的回忆。或许有一天，我会回想起这些时光，感慨万千；感谢敬爱的老师们，你们的教诲和指导让我受益匪浅。无论是在课堂上还是在课外，你们的耐心和专业精神都深深地影响了我，激励着我不断追求知识和成长。

感谢我的舍友们，与你们的相处很愉快，虽然我们在一起的时间不算长，但实在感谢你们的陪伴和支持。当我写下这句话时，已是深夜，所以要更加感谢你们的包容。

感谢我的朋友们，无论你们现在在哪里，将来又将去往何方。
在这段旅程中，你们给予了我无尽的支持和陪伴，让我的生活丰富多彩。其中，在大学认识的新朋友中，
我要尤其感谢武泽恺和罗剑波。与你们的深度交流和讨论，让我思考了许多，这份友谊是我大学生活中最宝贵的财富之一。

最后，我要感谢我的家人。你们的无私支持和鼓励是我能够坚持到现在的最大动力。无论是在学业上还是生活中，你们一直是我最坚强的后盾。感谢你们一直以来的理解和支持，我会继续努力，不辜负你们的期望。

沈从文写到：“凡事都有偶然的凑巧，结果却又如宿命的必然。”我时常想起这句话，回顾我的大学生活，确实有很多偶然的相遇和经历，但这些偶然的事件却又构成了我人生中必然的轨迹。正是这一个又一个偶然，让我成长为今天的自己。
我感恩这一路上的点点滴滴，无论是成功还是失败，无论是喜悦还是忧伤，无论是刻骨铭心的经历还是平凡的日常。
\end{acknowledgement}